%% kapa: exam=1 round=22 year=2011 period=2 form=A subject=appraisal_law subject_ko=부동산관계법규 seq=13 number=53
%% source: Q-Net 감정평가사 기출문제 (cst00309) artlSeq=5036790 file=extracted/1008/5036790/2011년 제22회 감정평가사 1차시험 기출문제/2교시.hwp
%% conversion: kordoc 4.24.51 -> tools/gen1.py (auto); review=proofread-gpt-astra; flags=none
\begin{question}{53}
국토의 계획 및 이용에 관한 법령상 도시계획시설사업 시행에 관한 설명으로 옳지 않은 것은?
\begin{choices}
  \item 특별시장ㆍ광역시장ㆍ시장 또는 군수는 도시계획시설에 대하여 도시계획시설 결정의 고시일부터 2년 이내에 대통령령으로 정하는 바에 따라 단계별 집행계획을 수립하여야 한다.
  \item 도시계획시설사업이 둘 이상의 특별시ㆍ광역시ㆍ시 또는 군의 관할 구역에 걸쳐 시행되게 되는 경우에는 관계 특별시장ㆍ광역시장ㆍ시장 또는 군수가 서로 협의하여 시행자를 정한다.
  \item 군수가 관할구역내에서 도시계획시설사업을 시행하는 경우에는 그 도시계획시설사업에 관한 실시계획을 작성한 후 국토해양부장관의 인가를 받아야 한다.
  \item 특별시장ㆍ광역시장ㆍ시장 또는 군수는 실시계획의 인가를 받지 아니하고 도시계획시설사업을 하거나 그 인가 내용과 다르게 도시계획시설사업을 하는 자에게 그 토지의 원상회복을 명할 수 있다.
  \item 국토해양부장관, 시ㆍ도지사 또는 대도시 시장은 실시계획을 작성한 경우에는 대통령령으로 정하는 바에 따라 그 내용을 고시하여야 한다.
\end{choices}
\end{question}
\answer{3}
